
\documentclass{article}
\usepackage{colortbl}
\usepackage{makecell}
\usepackage{multirow}
\usepackage{supertabular}

\begin{document}

\newcounter{utterance}

\centering \large Interaction Transcript for game `cladder', experiment `full\_v1.5\_default', episode 9125 with refineonly{-}Qwen3.5{-}27B.
\vspace{24pt}

{ \footnotesize  \setcounter{utterance}{1}
\setlength{\tabcolsep}{0pt}
\begin{supertabular}{c@{$\;$}|p{.15\linewidth}@{}p{.15\linewidth}p{.15\linewidth}p{.15\linewidth}p{.15\linewidth}p{.15\linewidth}}
    \# & \multicolumn{2}{c}{Player} && \multicolumn{2}{c}{Game Master} \\
    \hline

    \theutterance \stepcounter{utterance}  
    & & & \multicolumn{4}{p{0.6\linewidth}}{
        \cellcolor[rgb]{0.9,0.9,0.9}{
            \makecell[{{p{\linewidth}}}]{
                \texttt{\tiny{[P1$\langle$GM]}}
                \texttt{You are an expert in causal inference. The following question is not a typical commonsense query, but rather a meticulously designed question created by a professor specializing in causal inference, intended to assess the students' mastery of the course content.} \\
\texttt{Imagine a self{-}contained, hypothetical world with only the following conditions, and without any unmentioned factors or causal relationships: Kraz has a direct effect on pexu and rukz. Pexu has a direct effect on hwax. Hwax has a direct effect on rukz. Kraz is unobserved. For those who are not pexu and are not hwax, the probability of rukz is 19\%. For those who are not pexu and are hwax, the probability of rukz is 29\%. For those who are pexu and are not hwax, the probability of rukz is 20\%. For those who are pexu and are hwax, the probability of rukz is 31\%. For those who are not pexu, the probability of hwax is 16\%. For those who are pexu, the probability of hwax is 90\%. For those who are pexu, would it be less likely to see rukz if the individual was not pexu?} \\
\texttt{Start your answer with 'Yes' or 'No', followed by additional reasoning or evidence to support your explanation.} \\
            }
        }
    }
     \\ \\

    \theutterance \stepcounter{utterance}  
    & \multicolumn{4}{p{0.6\linewidth}}{
        \cellcolor[rgb]{0.561,0.737,0.561}{
            \makecell[{{p{\linewidth}}}]{
                \texttt{\tiny{[P1$\rangle$GM]}}
                \texttt{No} \\
\\ 
\texttt{**Reasoning:**} \\
\\ 
\texttt{To answer this question, we must compare the probability of the outcome *rukz* (\$R\$) for the subgroup of individuals who are *pexu* (\$P=1\$) under two scenarios: their actual state (\$P=1\$) and the counterfactual state where they are *not pexu* (\$P=0\$). The question asks if the probability of *rukz* would be **less** in the counterfactual scenario (\$P=0\$) compared to the factual scenario (\$P=1\$).} \\
\\ 
\texttt{**Step 1: Calculate the probability of *rukz* given *pexu* (\$P(R=1 | P=1)\$)**} \\
\texttt{We use the Law of Total Probability, conditioning on the mediator *hwax* (\$H\$).} \\
\texttt{Given:} \\
\texttt{*   \$P(R=1 | P=1, H=0) = 20\textbackslash{}\%\$} \\
\texttt{*   \$P(R=1 | P=1, H=1) = 31\textbackslash{}\%\$} \\
\texttt{*   \$P(H=1 | P=1) = 90\textbackslash{}\% \textbackslash{}implies P(H=0 | P=1) = 10\textbackslash{}\%\$} \\
\\ 
\texttt{Calculation:} \\
\texttt{\$\$P(R=1 | P=1) = P(R=1 | P=1, H=0) \textbackslash{}cdot P(H=0 | P=1) + P(R=1 | P=1, H=1) \textbackslash{}cdot P(H=1 | P=1)\$\$} \\
\texttt{\$\$P(R=1 | P=1) = (0.20 \textbackslash{}cdot 0.10) + (0.31 \textbackslash{}cdot 0.90)\$\$} \\
\texttt{\$\$P(R=1 | P=1) = 0.02 + 0.279 = 0.299 \textbackslash{}text\{ (or \} 29.9\textbackslash{}\%)\$\$} \\
\\ 
\texttt{**Step 2: Calculate the probability of *rukz* given *not pexu* (\$P(R=1 | P=0)\$)**} \\
\texttt{Similarly, we condition on *hwax* for the group where \$P=0\$.} \\
\texttt{Given:} \\
\texttt{*   \$P(R=1 | P=0, H=0) = 19\textbackslash{}\%\$} \\
\texttt{*   \$P(R=1 | P=0, H=1) = 29\textbackslash{}\%\$} \\
\texttt{*   \$P(H=1 | P=0) = 16\textbackslash{}\% \textbackslash{}implies P(H=0 | P=0) = 84\textbackslash{}\%\$} \\
\\ 
\texttt{Calculation:} \\
\texttt{\$\$P(R=1 | P=0) = P(R=1 | P=0, H=0) \textbackslash{}cdot P(H=0 | P=0) + P(R=1 | P=0} \\
            }
        }
    }
    & & \\ \\

    \theutterance \stepcounter{utterance}  
    & & & \multicolumn{2}{p{0.3\linewidth}}{
        \cellcolor[rgb]{0.9,0.9,0.9}{
            \makecell[{{p{\linewidth}}}]{
                \texttt{\tiny{[GM$|$GM]}}
                \texttt{No} \\
            }
        }
    }
    & & \\ \\

    \theutterance \stepcounter{utterance}  
    & & & \multicolumn{2}{p{0.3\linewidth}}{
        \cellcolor[rgb]{0.9,0.9,0.9}{
            \makecell[{{p{\linewidth}}}]{
                \texttt{\tiny{[GM$|$GM]}}
                \texttt{yes} \\
            }
        }
    }
    & & \\ \\

    \theutterance \stepcounter{utterance}  
    & & & \multicolumn{2}{p{0.3\linewidth}}{
        \cellcolor[rgb]{0.9,0.9,0.9}{
            \makecell[{{p{\linewidth}}}]{
                \texttt{\tiny{[GM$|$GM]}}
                \texttt{game\_result = LOSE} \\
            }
        }
    }
    & & \\ \\

\end{supertabular}
}

\end{document}
